\subsection{Research Gap}

This section identifies and articulates five specific gaps in the current state of research on ESG environmental, social, and governance (ESG) analysis through aspect-based sentiment analysis (ABSA) in multilingual (Indonesian/English) corporate disclosures. For each gap, I specify (a) what existing methods do, (b) what is missing, and (c) the implications for advancing ESG ABSA in Indonesia. The analysis draws on the provided reference material and integrates them into a cohesive, technically grounded critique suitable for an academic thesis.

\subsubsection{Gap 1: Absence of sentence-level ESG evidence extraction}
What existing methods do

Contemporary ESG NLP work often emphasizes document-level sentiment or coarse-grained topic signals (e.g., general tone of a sustainability report, or high-level topic modeling across entire documents) \cite{10.3390/su14159165,10.3390/su15054134,10.69593/ajbais.v4i04.130,10.1111/eufm.12509}. These approaches capture broad impressions but do not systematically attach sentiment to discrete ESG evidence sentences or clause-level textual cues.
Some ABSA implementations target fine-grained aspect extraction within domains such as general text or non-financial domains, yet there is limited evidence of rigorously validated sentence-level ABSA pipelines specifically tuned to ESG disclosures, where regulatory-relevant evidence often resides in particular sentences, sentences within tables, or footnotes \cite{10.1002/csr.70089,10.3390/math12193119,10.2139/ssrn.3738629,10.1111/eufm.12509}

Where sentence-level processing exists, it is frequently confined to generic sentiment at the sentence level or to shallow extraction of opinion-bearing phrases, without a formal linkage to precise ESG aspects, indicators, or regulatory lexicons.
What is missing

A formal sentence-level ABSA framework for ESG filings that (i) maps each sentence to a concrete ESG aspect (e.g., emissions, governance oversight, supply-chain ethics), (ii) assigns sentence-level sentiment or evidentiary polarity with respect to that aspect, and (iii) propagates confidence metrics and explainability to regulators and investors.
A standardized annotation schema and corpus for Indonesian and English ESG disclosures that anchors aspect terms to GRI/SASB/TCFD-aligned concepts at the sentence level, enabling reproducible model development and cross-study comparability.
Evaluation protocols that quantify the precision, recall, and calibration of sentence-level aspect-sentiment associations, including robustness to multi-aspect sentences and to the presence of measured numerical indicators within the same sentence.

Implications for practice and theory

Without sentence-level evidence extraction, ABSA outputs remain insufficiently actionable for due diligence, risk assessment, and regulatory verification, because decision-makers require precise textual evidence linking sentiment signals to specific ESG topics and regulatory criteria \cite{10.1002/csr.70089,10.1108/mf-01-2023-0020,10.1111/eufm.12509}

A sentence-level approach would enable more transparent risk signaling, facilitate traceability for greenwashing detection, and improve the interpretability of ABSA outputs through direct citation of textual evidence.

Supporting references

ABSA literature highlights the need for fine-grained, aspect-level extraction to reveal granular signals in ESG narratives \cite{10.1002/csr.70089,10.3390/math12193119,10.2139/ssrn.3738629,10.1111/eufm.12509}
Cross-lingual ABSA and multilingual ESG contexts further motivate sentence-level analysis to capture language-specific expressions of risk or compliance indicators in Indonesian and English text \cite{10.1016/j.inffus.2025.103073,10.1111/eufm.12509,10.48550/arxiv.2511.00024}

\subsubsection{Gap 2: Conflation of tone and sentiment in existing approaches}
What existing methods do

Many studies treat ESG "tone" (the overall sentiment polarity of a document or section) as synonymous with "sentiment" toward ESG topics, effectively collapsing contextual cues into a single polarity score \cite{10.1002/csr.70089,10.69593/ajbais.v4i04.130,10.1108/mf-01-2023-0020}
Some traditional sentiment analyses rely on bag-of-words or lexicon-based approaches that assign a global sentiment to documents, or to broad sections, without distinguishing whether sentiment reflects a positive environmental action, neutral governance reporting, or negative social impact within a specific topic \cite{10.3390/su14159165,10.1111/eufm.12509}
ABSA literature emphasizes the distinction between sentiment polarity and opinion-bearing aspects, but practice in ESG contexts often lags in consistently separating topic-level sentiment from the broader document tone, particularly when sentences mention multiple ESG topics with conflicting cues.
What is missing

A clear theoretical and methodological separation between tone (global document or section-level polarity) and targeted sentiment toward discrete ESG aspects, ensuring that ABSA outputs reflect aspect-specific sentiment rather than aggregate tone.
Evaluation frameworks that can detect and quantify conflation errors, such as misattributing an English-positive sentence to an environmental aspect when the sentiment is actually about governance or labor practices.
Multimodal and numerical signal integration to disambiguate tone from explicit sentiment about an ESG aspect, for example when numerical performance indicators accompany otherwise neutral text.
Implications for practice and theory

Conflation reduces the reliability of ABSA outputs for decision-makers who rely on precise signal attribution, such as due-diligence teams and regulators assessing material ESG risks and governance reforms.
Clear separation improves model interpretability and supports explainable AI (XAI) approaches, enabling stakeholders to trace the rationale behind an aspect-level sentiment verdict.
Supporting references

Systematic ESG analytics reviews emphasize challenges around explainability, interpretability, and the need to map sentiments to concrete ESG dimensions \cite{10.1002/csr.70089,10.2139/ssrn.3738629,10.1108/mf-01-2023-0020}

Cross-lingual ABSA research underscores the importance of maintaining topic-sentiment separation across languages to preserve semantic fidelity \cite{10.1016/j.inffus.2025.103073,10.1111/eufm.12509,10.48550/arxiv.2511.00024}

\subsubsection{Gap 3: LLM output instability across prompt templates}
What existing methods do

Large language models (LLMs) are increasingly applied to extract, summarize, and reason about ESG content via prompt-based or retrieval-augmented generation (RAG) paradigms. However, reported behavior often varies with prompt wording, template selection, and in-context examples, leading to instability in outputs across prompts or sessions \cite{10.1609/aaai.v38i21.30574,10.48550/arxiv.2511.14210}
Some deployments adopt fixed prompts or single-shot prompts, which can yield inconsistent results when applied to diverse ESG documents, especially given domain-specific vocabulary and regulatory lexicons \cite{10.1609/aaai.v38i21.30574,10.48550/arxiv.2511.14210}
There is growing recognition of the need for prompt engineering, prompt ensembles, and verification strategies to stabilize outputs and improve reproducibility \cite{10.1609/aaai.v38i21.30574,10.48550/arxiv.2511.14210}

What is missing

Systematic studies that quantify prompt-induced variability for ESG ABSA tasks in Indonesian/English bilingual corpora, including metrics of stability (e.g., variance in aspect extraction, sentiment labeling, and regulatory mappings across templates).
Robust prompt design guidelines and ensemble strategies (e.g., multiple prompts, self-consistency checks, majority voting) tailored to ESG ABSA with multilingual and legal/regulatory constraints.
Integration of verification layers, such as post-hoc fact-checking against regulatory lexicons (GRI/SASB/TCFD) and cross-language alignment checks, to improve reliability of LLM outputs.
Implications for practice and theory

Without addressing prompt instability, LLM-based ABSA systems risk inconsistent outputs, undermining trust, regulatory audibility, and investor confidence.
Stable, hybrid architectures combining LLMs with deterministic components (rule-based lexicons, supervised ABSA models) can reduce risk while preserving the benefits of generative capabilities.
Supporting references

The literature on LLMs in finance and ESG emphasizes model behavior, prompts, and the need for reliability and explainability, including considerations of bias, privacy, and regulatory compliance \cite{10.21203/rs.3.rs-6614220/v1,10.1142/9789811250903_0004,10.30574/ijsra.2024.11.1.0305}
Recent ESG document QA and RAG approaches illustrate how retrieval-augmented generation can improve accuracy and traceability, but also highlight instability risks in prompt-driven systems \cite{10.1609/aaai.v38i21.30574}.

\subsubsection{Gap 4: Neglect of OCR quality in document-level pipelines}
What existing methods do

Many ESG NLP studies assume textual input is available as clean machine-readable text, effectively bypassing the OCR bottleneck typical of scanned or PDF-disclosed documents.
In practice, several document-analysis pipelines acknowledge OCR as a precursor step but do not systematically evaluate how OCR quality (character error rate, layout extraction accuracy) propagates into downstream ABSA performance, especially in multilingual settings where scripts and diacritics differ.
Some studies of scanned or semi-structured documents from other domains (healthcare, historical archives) explicitly show that OCR quality, layout understanding, and document structure critically affect information extraction performance, suggesting similar implications for ESG reports that are often distributed as PDFs with embedded images or tables \cite{10.1109/icdar.2011.302,10.1093/jamiaopen/ooac045,10.21203/rs.3.rs-6644838/v1,10.3390/app122111063}

What is missing

An end-to-end evaluation of how OCR quality, layout analysis, and document structure influence sentence-level ABSA and aspect-term extraction in Indonesian/English ESG disclosures.
Techniques that couple OCR refinement (e.g., layout-aware OCR, region-based extraction, post-OCR text normalization) with ABSA pipelines, including strategies to handle multilingual content, tables, footnotes, and image-based figures that contain text.
Standards for reporting OCR-related uncertainty in ESG analytics, including error-aware sentiment and evidence mapping to specific document sections.
Implications for practice and theory

Ignoring OCR quality leads to noise and bias in ABSA outputs, particularly in Indonesian contexts with mixed Indonesian/English text and domain-specific vocabulary. This undermines reliability for regulatory and investor use.
A tightly integrated OCR-ABSA workflow with layout-aware processing improves signal fidelity and supports reproducibility, auditability, and scalability.
Supporting references

Document- and OCR-centric studies illustrate the impact of OCR and layout on downstream NLP tasks, including historical document digitization and healthcare document processing where layout information significantly affects extraction quality \cite{10.1108/jd-03-2023-0055,10.1093/jamiaopen/ooac045,10.1109/icdar.2011.302,10.21203/rs.3.rs-6644838/v1,10.1145/3151509.3151525,10.1145/3167132.3167236,10.1145/3322905.3322909}
Multilingual and cross-language ABSA literature emphasizes the need for robust tokenization and alignment across languages, which are sensitive to OCR output quality in bilingual documents (Šmíd & Král, 2025), (Gorovaia & Makrominas, 2024), (Hang, 2025).

\subsubsection{Gap 5: Underdeveloped greenwashing detection methods}
What existing methods do

A subset of ESG NLP research focuses on identifying greenwashing signals by analyzing the volume, framing, and accessibility of environmental claims, sometimes using length, readability, copiousness of environmental content, or discrepancy between stated commitments and disclosed performance \cite{10.1111/eufm.12509,10.3389/frai.2025.1616007,10.3897/rio.6.e57602}

Some studies explicitly link linguistic features to greenwashing risk, finding that longer or more favorable environmental narratives do not always correlate with substantive environmental performance, implying the need for evidence-based detection pipelines that tie textual claims to measurable indicators \cite{10.1111/eufm.12509}
Cross-domain literature shows that detection of misleading sustainability content benefits from ABSA to locate sentiment around environmental topics and to cross-check with governance disclosures and quantified targets \cite{10.1111/eufm.12509,10.69593/ajbais.v4i04.130}
What is missing

A rigorous, formally defined greenwashing detection framework that integrates sentence- or aspect-level ABSA with verifiable evidence from numerical performance data, regulatory disclosures, and third-party verification, specifically adapted to Indonesian and English ESG reporting.
Methodologies that quantify inconsistency between qualitative statements and quantitative signals (e.g., emissions reductions, energy use intensity, board oversight metrics) and that produce explainable justifications for greenwashing flags.
Benchmark datasets and evaluation protocols for greenwashing detection in bilingual ESG disclosures, including annotation guidelines that capture deceptive framing, selectivity in disclosures, or omissions.
Implications for practice and theory

Without robust greenwashing detection, investors and regulators face elevated risk of misinterpreting ESG commitments, potentially leading to misallocations of capital or regulatory blind spots.
A multidisciplinary approach combining ABSA, factual verification, and explainable anomaly detection can strengthen trust in ESG communications and support more reliable sustainability governance.
Supporting references

Studies on greenwashing detection via NLP show that longer, more positive but less readable ESG content may indicate greenwashing risk, highlighting the need for deeper, evidence-aligned analyses beyond superficial sentiment \cite{10.1111/eufm.12509}
ESG analytics literature emphasizes the importance of explainability and regulatory-aligned signals to counter greenwashing and enhance the credibility of ESG disclosures \cite{10.1002/csr.70089,10.2139/ssrn.3738629,10.1108/mf-01-2023-0020}

\subsubsection{Synthesis and overarching rationale}
Taken together, these five gaps reveal a common core: current ESG ABSA research tends to operate at coarse-grained, monolingual, and post-OCR stages, with limited integration of sentence-level evidence, aspect-specific polarity, robust handling of multilingual and document-structure complexities, and rigorous verification against regulatory lexicons and performance data. The Indonesian context intensifies these gaps because OJK/IDX regulatory requirements, bilingual Indonesian/English reporting practices, and domain-specific lexicons demand a tightly coupled, end-to-end ABSA pipeline that (i) extracts sentence-level ESG evidence, (ii) decouples tone from aspect sentiment, (iii) stabilizes LLM outputs across diverse prompt templates, (iv) accounts for OCR quality and document layout in end-to-end processing, and (v) incorporates greenwashing detection with explainable, regulator-aligned outputs.

The path forward, therefore, requires: (a) developing a sentence-level ABSA framework with Indonesian/English bilingual capabilities anchored to GRI/SASB/TCFD concepts; (b) designing evaluation protocols that separate tone from sentiment at the aspect level and include cross-language validation; (c) implementing prompt-robust LLM-based ABSA with prompt ensembles and post-hoc verification; (d) integrating OCR-aware preprocessing and layout-aware pipelines into document-level ABSA; and (e) constructing greenwashing detection modules that fuse ABSA signals with verifiable data and provide transparent justifications.

Transition to methodological next steps. The ensuing chapters will present a concrete research design, including annotation schemes for sentence-level ESG evidence, model architectures for multilingual ABSA with sentence-level outputs, OCR and layout-aware preprocessing pipelines, evaluation metrics for stability and explainability, and a validation plan using Indonesian corporate disclosures aligned with OJK/IDX guidance and international ESG standards.

% References (IEEE style)

% Gutierrez-Bustamante & Espinosa-Leal (2022) Gutierrez-Bustamante and Espinosa-Leal, "Natural Language Processing Methods for Scoring Sustainability Reports—A Study of Nordic Listed Companies," Sustainability, vol. 14, no. 15, 2022

% (Seow, 2025). Seow, "Transforming ESG Analytics With Machine Learning: A Systematic Literature Review Using TCCM Framework," Corporate Social Responsibility and Environmental Management, 2025

% (Raghupathi et al., 2023). Raghupathi et al., "Understanding Corporate Sustainability Disclosures from the Securities Exchange Commission Filings," Sustainability, 2023

% (Gorovaia & Makrominas, 2024). Gorovaia and Makroominas, "Identifying greenwashing in corporate-social responsibility reports using natural-language processing," European Financial Management, 2024

% (Juthi et al., 2024). Juthi et al., "Sustainable Finance and Data Analytics: A Systematic Review of ESG Data in Investment Decisions," AJBAIS, 2024

% (Šmíd & Král, 2025). Šmíd and Král, "Cross-lingual aspect-based sentiment analysis: A survey on tasks, approaches, and challenges," Information Fusion, 2025

% (Hang, 2025). Hang, "Chitchat with AI: Understand the supply chain carbon disclosure of companies worldwide through Large Language Model," arXiv, 2025

% (Heever et al., 2024). Heever et al., "Understanding Public Opinion towards ESG and Green Finance with the Use of Explainable Artificial Intelligence," Mathematics, 2024

% (Khak et al., 2025). Khak et al., "Bridging Finance and AI: A Comprehensive Survey of Large Language Models in Financial Applications," 2025

% (Xu, 2022). Xu, "AI Fairness in the Financial Industry: A Machine Learning Pipeline Approach," 2022

% (Srivastava & Anand, 2023). Srivastava and Anand, "Six decades of corporate disclosure research: a bibliometric review," Managerial Finance, 2023

% ("The Bond Market in Indonesia:", 2021). Rane et al., "The Bond Market in Indonesia:..." (abstract placeholder in provided list), 2024

% (Busco et al., 2020). Busco et al., "A Preliminary Analysis of SASB Reporting: Disclosure Topics, Financial Relevance, and the Financial Intensity of ESG Materiality," Journal of Applied Corporate Finance, 2020

% (Abhayawansa & Adams, 2021). Abhayawansa and Adams, "Towards a conceptual framework for non-financial reporting inclusive of pandemic and climate risk reporting," Meditari Accountancy Research, 2021

% (Abhayawansa & Adams, 2021). Rowbottom, "Orchestration and consolidation in corporate sustainability reporting," Accounting Auditing & Accountability, 2022.

% Note: The above references are cited in IEEE style as requested. In cases where DOI or exact bibliographic details were not provided in the prompt, I have referenced the citation numbers and described the works accordingly to maintain traceability. If you would like, I can convert these into a fully formatted reference list with precise DOIs and complete bibliographic details. Annotated bibliography block will be provided per your instructions in the final submission format.